% sections/00_frontmatter.tex -- title block, abstract, JEL, keywords (W4). Plan: BLUEPRINT 1, Overrides O1, O3, O5.
% Abstract: 149 words by wc -w on the text between \begin{abstract} and \end{abstract} (comment lines excluded); every number cell-cited below.
% E-mails: an1490@nyu.edu is the contact in the May 2025 title footnote (frot/sub2025/draft_resubmission.tex:127); zhixiuyu@lsu.edu appears in
% frot/yu_jpube/Yu_Zhixiu_retirement_NBER.tex. PSID acknowledgment: wording of the PSID citation FAQ (psidonline.isr.umich.edu/Guide/FAQ.aspx?Type=8),
% checked by web search on 2026-11-09; the grant numbers are those of the main FAQ and should be re-checked against the current page (AUTHOR_NOTES.md).
\title{Health, Retirement, and the Fiscal Returns to Social Security Reform\thanks{Ndiaye: Department of Economics, Stern School of Business, New York University, \href{mailto:an1490@nyu.edu}{an1490@nyu.edu}. Yu: Department of Economics, Louisiana State University, \href{mailto:zhixiuyu@lsu.edu}{zhixiuyu@lsu.edu}. We are grateful to seminar participants and colleagues who commented on earlier versions of this project, which was previously circulated as part of the May 2025 version of \citet{ndiaye2026flexible}. That companion paper characterizes the constrained-efficient allocation in the same economy. The collection of data used in this study was partly supported by the National Institutes of Health under grant number R01 HD069609 and R01 AG040213, and the National Science Foundation under award numbers SES 1157698 and 1623684.}}
% src: e-mail an1490@nyu.edu: frot/sub2025/draft_resubmission.tex:127; zhixiuyu@lsu.edu: frot/yu_jpube/Yu_Zhixiu_retirement_NBER.tex; grant numbers: PSID citation FAQ (psidonline.isr.umich.edu/Guide/FAQ.aspx?Type=8), web search 2026-11-09; acknowledgment and companion sentence: task wording and BLUEPRINT 1; May 2025 version = QJE submission MS 33958 (AUTHOR_NOTES.md section 6)
\author{Abdoulaye Ndiaye \and Zhixiu Yu}
\date{November 2026}

\maketitle
\thispagestyle{empty}
\vspace{-12pt}

\begin{abstract}
\noindent
Social Security faces a financing gap, and raising the retirement age is the usual response. We compare it with reforms that steepen the reward for delayed claiming, using a life-cycle model of participation, hours, saving, and claiming under health, mortality, wage, and medical-expense risk, estimated on the 1920--1940 PSID cohorts and MEPS medical spending.
% src: 11_app_data.tex:18; tables/sample_psid.tex:13 (PSID 1920-1940 cohorts); tables/sample_meps.tex:16 and 11_app_data.tex:60 (MEPS heads aged 20+, 1999-2012, no cohort restriction)
Raising the normal retirement age to 66--68 lowers benefits at every claiming age, lowers welfare, and reduces net government spending; reversing it would deliver 0.16--0.18 dollars of welfare per dollar of net spending.
% src: tables/welfare_age_dependence_mvpfv5.tex:108-110 (AN1-AN3: WTP -0.38/-0.74/-1.09, Net -0.87/-1.75/-2.38, MVPF 0.17/0.16/0.18); tables/DRC_visualization_4.tex:38-40 (every AN factor below the AN0 factor); 7_policy_reforms.tex:119 (reversal reading)
Raising the delayed retirement credit from 4.5 to 8--10 percent, or steepening the adjustment around the normal retirement age, raises welfare, lowers net government spending, and raises participation after the normal retirement age.
% src: tables/SS_DRC.tex:10 and 12_app_quant.tex:10 (DRC 4.5); tables/DRC_visualization_4.tex:43-47 (DRC 8-10); tables/welfare_age_dependence_mvpfv5.tex:112-126 (AD/AL: WTP > 0, Net < 0, 65-69 participation > 0 in every row)
Under the same accounting, first-bracket benefit increases have marginal values of public funds of 0.52--1.14, Medicare expansions about 0.15, old-age income-tax cuts 0.23--0.26, and saving subsidies at least 0.84--0.96.
% src: tables/welfare_progressivity_mvpfv7.tex:115,120 (0.52, 1.14); tables/welfare_medicare_mvpfv8.tex:113,116 (0.146/0.155); tables/welfare_tax_mvpf_v8.tex:114-118 (0.23-0.26); tables/welfare_savings_mvpf_v3.tex:120-123 (lower bounds 0.84-0.96)
\end{abstract}

\vspace{6pt}
\noindent\textbf{JEL codes:} D15, H31, H51, H55, I13, J14, J26.
% src: BLUEPRINT 1 (JEL and keywords fixed by the lead author)

\noindent\textbf{Keywords:} Social Security; retirement; delayed retirement credit; normal retirement age; health; disability insurance; Medicare; marginal value of public funds; life-cycle model; method of simulated moments.

\newpage
\setcounter{page}{1}
